\documentclass[../../main.tex]{subfiles}
\begin{document}
\begin{oframed}
  \subfile{problem}
\end{oframed}

{\bf [解答]}
与えられた関数方程式
\begin{align}
  f(x+y)=f(x)+f(y)+f(x)f(y) \label{eq:1}
\end{align}
について考える．

\paragraph{(1)}

まず\cref{eq:1}で$(x,y)=(0,0)$として，
\begin{align}
  f(0)(1+f(0))=0 \quad\therefore\ f(0)=0,-1 \label{eq:2}
\end{align}
である．次に$(x,y)=(x,0)$とすると
\begin{align}
  & f(x) = f(x)+f(0)+f(x)f(0) \\
  & \left\{f(x)+1\right\}f(0) = 0 \\
  \therefore
  & f(0) = 0 \lor f(x) \equiv -1 \label{eq:3}
\end{align}
である．\cref{eq:2,eq:3}より，$f(0)=0$または恒等的に$f(x)=-1$かのどちらかである．

さて\cref{eq:1}から
\begin{align}
  f(x+y)-f(x) = f(y)\left(1+f(x)\right)
\end{align}
だから，$y \neq 0$として両辺$y$で割ると
\begin{align}
  \dfrac{f(x+y)-f(x)}{y}=\dfrac{f(y)(1+f(x))}{y}
\end{align}
となる．両辺$y\to 0$の極限をとると
\begin{align}
  \lim_{y\to 0} \dfrac{f(x+y)-f(x)}{y} = \lim_{y\to 0}\dfrac{f(y)}{y}(1+f(x))
\end{align}
である．ここで右辺の極限について\cref{eq:2}および題意より$f'(0)$が存在することより
\begin{align}
  \lim_{y\to 0}\dfrac{f(y)}{y}(1+f(x)) =
  \begin{dcases}
    \lim_{y\to 0}\dfrac{f(y)-f(0)}{y}(1+f(x)) =  f'(0)(1+f(x)) & (f(0)=0) \\
    \lim_{y\to 0}\dfrac{f(y)-f(0)}{y}(1+f(x)) = 0 & (f(x)=-1)
  \end{dcases}
\end{align}
だから代入して
\begin{align}
  \lim_{y\to 0} \dfrac{f(x+y)-f(x)}{y} = f'(0)(1+f(x)) \label{eq:4}
\end{align}
である．ただし$f(x)=-1$の場合も$1+f(x)=0$よりこれで良いことを用いた．
従って左辺の極限値が存在するから，$f(x)$は任意の$x$で微分可能である．

\paragraph{(2)}

\cref{eq:4}より
\begin{align}
  f'(x) = f'(0)(1+f(x))
\end{align}
だから$y=f(x)$として
\begin{align}
  \dfrac{dy}{dx}=f'(0)(1+y)
\end{align}
となる．まず$y=-1$はこの式を満たし，\cref{eq:1}も満たすから十分で解である．

そこで以下$y\neq -1$とする．両辺$1+y$で割って
\begin{align}
  \dfrac{dy}{1+y} = f'(0) dx \label{eq:5}
\end{align}
である．両辺を積分して
\begin{align}
  & \log |1+y| =f'(0)x+C \\
  \therefore
  & y=e^C \cdot e^{f'(0)x}-1 \label{eq:6}
\end{align}
とおける．ただし$C$は積分定数である．
\cref{eq:2}に代入して$C=0$である．

この時\cref{eq:6}より$y=e^{f'(0)x}-1$だから，\cref{eq:1}に代入すると
\begin{align}
  \begin{split}
    (\text{左辺})&=e^{f'(0)(x+y)}-1 \\
    (\text{右辺})&=e^{f'(0)x}+e^{f'(0)y}-2+(e^{f'(0)x}-1)(e^{f'(0)y}-1)=e^{f'(0)(x+y)}-1
  \end{split}
\end{align}
となって恒等的に\cref{eq:1}が成立するから十分である．

以上から求める$f(x)$は
\begin{align}
  y=-1,\quad e^{f'(0)x}-1
\end{align}
である．

\newpage

\end{document}
